\section{Continuidad}\label{sec:cont}
\subsection{De $\varepsilon$--$\delta$ a los abiertos}
En cálculo, $f:\R\to\R$ es continua en $p$ si para todo $\varepsilon>0$ existe $\delta>0$ tal que
$|x-p|<\delta\Rightarrow|f(x)-f(p)|<\varepsilon$. En palabras: \emph{para todo intervalo $V$ alrededor de $f(p)$ existe
un intervalo $U$ alrededor de $p$ con $f(U)\subseteq V$}. Si se cambia «intervalo alrededor de» por «abierto que
contiene a», la frase ya no menciona distancias y tiene sentido en cualquier espacio topológico.

\begin{definicion}[Definición: continuidad en un punto y continuidad global]
Sea $f:(X,\tau_X)\to(Y,\tau_Y)$.
\begin{itemize}
\item $f$ es \textbf{continua en $p\in X$} si para todo abierto $V$ de $Y$ con $f(p)\in V$ existe un abierto $U$ de $X$
con $p\in U$ y $f(U)\subseteq V$.
\item $f$ es \textbf{continua} si la imagen inversa de todo abierto de $Y$ es un abierto de $X$:
\[V\in\tau_Y\ \Longrightarrow\ f^{-1}(V)\in\tau_X .\]
\end{itemize}
\end{definicion}
Recuérdese que $f^{-1}(V)=\{x\in X: f(x)\in V\}$ está definido para \emph{cualquier} función; no supone que $f$ tenga
inversa.

\begin{teorema}[Equivalencias de la continuidad]
Para $f:X\to Y$ son equivalentes:
\begin{enumerate}[label=(\alph*),leftmargin=1.9em]
\item $f$ es continua (imagen inversa de abierto es abierto);
\item la imagen inversa de todo \emph{cerrado} de $Y$ es cerrada en $X$;
\item la imagen inversa de todo miembro de una base (o de una sub-base) de $Y$ es abierta en $X$;
\item $f$ es continua en cada punto de $X$;
\item $f(\cl{A})\subseteq\cl{f(A)}$ para todo $A\subseteq X$.
\end{enumerate}
\end{teorema}
\emph{Demostración de (a)$\iff$(b).} $f^{-1}(Y\sm V)=X\sm f^{-1}(V)$. \qed

\emph{Demostración de (a)$\iff$(d).} ($\Rightarrow$) Dado $p$ y $V\ni f(p)$ abierto, $U=f^{-1}(V)$ es abierto, contiene
a $p$ y $f(U)\subseteq V$. ($\Leftarrow$) Sea $V$ abierto de $Y$ y $p\in f^{-1}(V)$. Por la continuidad en $p$ hay un
abierto $U_p\ni p$ con $f(U_p)\subseteq V$, es decir, $U_p\subseteq f^{-1}(V)$. Entonces $f^{-1}(V)$ es vecindad de cada
uno de sus puntos y, por (N4), es abierto. \qed

\begin{metodo}[Método: decidir si $f$ es continua (espacios finitos)]
\begin{enumerate}[leftmargin=1.6em]
\item Escribir la tabla de valores de $f$ (leerla del diagrama de flechas con cuidado).
\item Para \emph{cada} abierto $V$ de la llegada, calcular $f^{-1}(V)$: los puntos de la salida cuya flecha cae en $V$.
\item Si todas las imágenes inversas son abiertos de la salida, $f$ es continua. Si alguna no lo es, $f$ no es
continua y ese $V$ es el contraejemplo.
\end{enumerate}
Para $\vac$ y para el total no hay nada que revisar: $f^{-1}(\vac)=\vac$ y $f^{-1}(Y)=X$.
\end{metodo}
\begin{metodo}[Método: decidir si $f$ es continua en un punto $p$]
\begin{enumerate}[leftmargin=1.6em]
\item Calcular $f(p)$ y listar los abiertos $V$ de la llegada que contienen a $f(p)$.
\item Para cada uno, ver si $f^{-1}(V)$ es vecindad de $p$ (contiene un abierto que contiene a $p$).
\item En espacios finitos basta una sola prueba: $f$ es continua en $p$ $\iff f(U_p)\subseteq V_{f(p)}$, donde $U_p$ y
$V_{f(p)}$ son los abiertos mínimos de $p$ y de $f(p)$.
\end{enumerate}
\end{metodo}

\begin{ejemplo}[Ejemplo: problema IV.1]
$X=\s{a,b,c,d}$, $\tau=\{X,\vac,\s{a},\s{b},\s{a,b},\s{b,c,d}\}$ y $f(a)=b$, $f(b)=d$, $f(c)=b$, $f(d)=c$.

\emph{Global.} $f^{-1}(\s{b})=\s{a,c}$ (los puntos que van a $b$ son $a$ y $c$), y $\s{a,c}\notin\tau$. No es continua.

\emph{En $d$.} $f(d)=c$; abiertos que contienen a $c$: $\s{b,c,d}$ y $X$. $f^{-1}(\s{b,c,d})=X$, que es vecindad de
$d$. Sí es continua en $d$. Con abiertos mínimos: $U_d=\s{b,c,d}$, $f(U_d)=\s{d,b,c}\subseteq U_{f(d)}=U_c=\s{b,c,d}$.

\emph{En $c$.} $f(c)=b$; $U_c=\s{b,c,d}$, $f(U_c)=\s{b,c,d}\not\subseteq U_b=\s{b}$. No es continua en $c$.

La función es continua en $a$, $b$ y $d$, y discontinua solo en $c$: basta un punto de discontinuidad para que la
función no sea continua.
\end{ejemplo}

\begin{ejemplo}[Ejemplos generales que conviene recordar]
\begin{itemize}
\item Toda función constante es continua (la imagen inversa de un abierto es $\vac$ o $X$).
\item La identidad $(X,\tau)\to(X,\tau)$ es continua.
\item Si $X$ tiene la topología \textbf{discreta}, toda función $X\to Y$ es continua.
\item Si $Y$ tiene la topología \textbf{indiscreta}, toda función $X\to Y$ es continua.
\item La composición de funciones continuas es continua: $(g\circ f)^{-1}(W)=f^{-1}(g^{-1}(W))$.
\item La restricción de una función continua a un subespacio es continua.
\end{itemize}
\end{ejemplo}

\begin{ojo}[Errores frecuentes en continuidad]
\begin{itemize}
\item Revisar imágenes directas $f(U)$ en lugar de imágenes inversas $f^{-1}(V)$. Eso es la definición de función
\emph{abierta}, que es otra cosa.
\item Usar los abiertos del espacio equivocado: las $V$ se toman en la \emph{llegada} y se comprueba que $f^{-1}(V)$
esté en la topología de la \emph{salida}.
\item Leer mal el diagrama: $f^{-1}(\s{b})$ son todos los puntos de los que \emph{sale} una flecha que llega a $b$.
\item Creer que «continua en algunos puntos» implica continua.
\end{itemize}
\end{ojo}

% =====================================================================
\section{Funciones abiertas y funciones cerradas}
\begin{definicion}[Definición]
$f:X\to Y$ es \textbf{abierta} si $f(U)$ es abierto en $Y$ para todo abierto $U$ de $X$; es \textbf{cerrada} si
$f(F)$ es cerrado en $Y$ para todo cerrado $F$ de $X$.
\end{definicion}
Continua, abierta y cerrada son tres propiedades \textbf{independientes}: una función puede tener cualquiera de las
ocho combinaciones. La razón de fondo es que la imagen inversa se comporta bien con todas las operaciones de conjuntos,
mientras que la imagen directa solo respeta las uniones:
\begin{gather*}
f(A\cup B)=f(A)\cup f(B),\qquad f(A\cap B)\subseteq f(A)\cap f(B),\\
f(X\sm A)\ \text{y}\ Y\sm f(A)\ \text{no se relacionan en general}.
\end{gather*}
Por eso «abierta» no equivale a «cerrada», salvo en el caso biyectivo.

\begin{teorema}[Hechos útiles]
\begin{enumerate}[label=(\alph*),leftmargin=1.9em]
\item Para saber si $f$ es abierta basta revisar las imágenes de los miembros de una \emph{base} de $X$ (porque
$f(\bigcup_iB_i)=\bigcup_if(B_i)$).
\item Si $f$ es abierta, entonces $f(X)$ es abierto en $Y$; si es cerrada, $f(X)$ es cerrado en $Y$. (Dos pruebas
rápidas para descartar.)
\item Si $f$ es \textbf{biyectiva}: $f$ abierta $\iff f$ cerrada $\iff f^{-1}$ continua.
\end{enumerate}
\end{teorema}
\emph{Demostración de (c).} Para una biyección, $f(X\sm A)=Y\sm f(A)$, lo que da la primera equivalencia. Y la imagen
inversa de $U$ por la función $f^{-1}:Y\to X$ es precisamente $f(U)$; así, «$f^{-1}$ continua» dice «$f(U)$ abierto para
todo abierto $U$». \qed

\begin{metodo}[Método: decidir si $f$ es abierta / cerrada]
\textbf{Abierta:} para cada abierto $U$ de la salida calcular $f(U)$ (adónde llegan las flechas que salen de $U$) y
ver si está en la topología de la llegada.\\[2pt]
\textbf{Cerrada:} listar los cerrados de la salida y los de la llegada; para cada cerrado $F$ de la salida calcular
$f(F)$ y ver si está en la lista de cerrados de la llegada.\\[2pt]
\textbf{Descartes rápidos:} mirar $f(X)$; mirar los abiertos y cerrados de un solo punto.
\end{metodo}

\begin{ejemplo}[Ejemplo: problema IV.2, $f(x)=x^2$ en $\R$]
$f$ es continua y \textbf{no es abierta}: $f((-1,1))=[0,1)$, que no es abierto porque ningún intervalo alrededor de $0$
cabe en $[0,1)$\fuente{20}. Geométricamente: la parábola «dobla» la recta en $0$, y el punto del doblez queda en el
borde de la imagen.
\end{ejemplo}

\begin{ejemplo}[Ejemplo: problema IV.3, resumen]
Con $X=\s{1,2,3,4}$, $Y=\s{a,b,c,d}$, $\tau_2=\{Y,\vac,\s{c},\s{c,d},\s{a,b,c}\}$ y la $\tau_1$ del enunciado:
\begin{center}\footnotesize\renewcommand{\arraystretch}{1.25}\setlength{\tabcolsep}{4pt}
\begin{tabular}{@{}lcccl@{}}\toprule
función & continua & abierta & cerrada & por qué\\ \midrule
$f=(1\!\mapsto\! a,2\!\mapsto\! a,3\!\mapsto\! d,4\!\mapsto\! c)$ & \no & \no & \no & $f^{-1}(\s{c})=\s{4}$; $f(\s{2})=\s{a}$; $f(\s{1})=\s{a}$\\
$g=(1\!\mapsto\! b,2\!\mapsto\! c,3\!\mapsto\! c,4\!\mapsto\! d)$ & \si & \no & \no & $g(X)=\s{b,c,d}$: ni abierto ni cerrado\\
$h=(a\!\mapsto\! 2,b\!\mapsto\! 3,c\!\mapsto\! 4,d\!\mapsto\! 1)$ & \no & \no & \no & $h^{-1}(\s{2})=\s{a}$; $h(\s{c})=\s{4}$\\
$k=(a\!\mapsto\! 2,b\!\mapsto\! 4,c\!\mapsto\! 3,d\!\mapsto\! 3)$ & \no & \si & \no & $k^{-1}(\s{2})=\s{a}$; $k(Y)=\s{2,3,4}$ no cerrado\\ \bottomrule
\end{tabular}\end{center}
$g$ es continua y no abierta; $k$ es abierta y no continua: el mismo problema contiene los dos contraejemplos que
muestran la independencia de ambas nociones. Nótese el uso del descarte (b): $g(X)$ y $k(Y)$.
\end{ejemplo}

% =====================================================================
\section{Homeomorfismos}
\subsection{Definición e idea}
Dos espacios son «el mismo espacio con los puntos renombrados» cuando hay una correspondencia entre sus puntos que
también hace corresponder sus abiertos. Eso es un homeomorfismo: el isomorfismo de la topología.
\begin{definicion}[Definición: homeomorfismo]
$f:X\to Y$ es un \textbf{homeomorfismo} si es biyectiva, continua y su inversa $f^{-1}:Y\to X$ es continua. Si existe
un homeomorfismo se dice que $X$ y $Y$ son \textbf{homeomorfos} (o topológicamente equivalentes): $X\homeo Y$.
\end{definicion}
\begin{teorema}[Formas equivalentes]
Para una \emph{biyección} $f:X\to Y$ son equivalentes:
\begin{enumerate}[label=(\alph*),leftmargin=1.9em]
\item $f$ es homeomorfismo;
\item $f$ es continua y abierta;
\item $f$ es continua y cerrada;
\item $U$ es abierto en $X$ $\iff$ $f(U)$ es abierto en $Y$; es decir, $U\mapsto f(U)$ es una biyección de $\tau_X$
sobre $\tau_Y$.
\end{enumerate}
\end{teorema}
\begin{teorema}[$\homeo$ es una relación de equivalencia]
Reflexiva: la identidad es homeomorfismo. Simétrica: si $f$ es homeomorfismo, $f^{-1}$ también. Transitiva: la
composición de homeomorfismos es homeomorfismo.
\end{teorema}

\begin{ojo}[Biyectiva y continua no basta]
La inversa debe ser continua también. Ejemplo finito: la identidad de $X=\s{a,b}$ con la topología discreta hacia
$X$ con la indiscreta es biyectiva y continua, pero su inversa no es continua. Ejemplo clásico:
$t\mapsto(\cos t,\sen t)$ de $[0,2\pi)$ sobre la circunferencia es biyectiva y continua, y no es homeomorfismo
(la inversa «rompe» la circunferencia en el punto $(1,0)$)\fuente{22,23}.
\end{ojo}

\subsection{Invariantes topológicos: cómo probar que \emph{no} son homeomorfos}
Una propiedad es un \textbf{invariante topológico} si, cuando la tiene un espacio, la tienen todos los homeomorfos a
él. Si $X$ tiene un invariante y $Y$ no, entonces $X\not\homeo Y$. Para espacios finitos los más útiles son:
\begin{center}\small\renewcommand{\arraystretch}{1.25}
\begin{tabular}{@{}p{6.4cm}p{9.0cm}@{}}\toprule
\textbf{Invariante} & \textbf{Por qué se conserva}\\ \midrule
Número de puntos & Un homeomorfismo es una biyección.\\
Número de abiertos (y de cerrados) & $U\mapsto f(U)$ es biyección entre las topologías.\\
Número de abiertos de cada tamaño & $f(U)$ tiene tantos puntos como $U$.\\
Número de puntos abiertos; de puntos cerrados & Caso particular del anterior.\\
Número de conjuntos abiertos y cerrados a la vez & $f$ conserva abiertos y cerrados.\\
Existencia de dos abiertos no vacíos disjuntos & $f(U\cap V)=f(U)\cap f(V)$ para biyecciones.\\
La forma del orden de inclusión entre abiertos & $U\subseteq V\iff f(U)\subseteq f(V)$.\\ \bottomrule
\end{tabular}\end{center}

\begin{ojo}[El mismo número de abiertos no garantiza homeomorfismo]
En $X=\s{a,b,c}$, $\tau=\{\vac,\s{a},\s{a,b},X\}$ y $\tau'=\{\vac,\s{a},\s{b,c},X\}$ tienen 4 abiertos cada una y no son
homeomorfas: $\tau'$ tiene un abierto propio que también es cerrado ($\s{a}$) y $\tau$ no; o bien, $\tau'$ tiene dos
abiertos no vacíos disjuntos y $\tau$ no. Los invariantes sirven para \emph{negar}; para \emph{afirmar} hay que
construir el homeomorfismo.
\end{ojo}

\begin{metodo}[Método: decidir si $X\homeo Y$ (espacios finitos)]
\begin{enumerate}[leftmargin=1.6em]
\item \textbf{Descartar:} comparar número de puntos, número de abiertos y tamaños de los abiertos. Si algo difiere,
no son homeomorfos y ese es el argumento.
\item \textbf{Construir:} si todo coincide, emparejar los abiertos «por tamaño y posición»: los puntos abiertos con
los puntos abiertos, los pares abiertos con los pares abiertos, etc. Eso dicta adónde debe ir cada punto.
\item \textbf{Verificar:} escribir la tabla $U\mapsto\varphi(U)$ para todos los abiertos de $X$ y comprobar que la
fila de imágenes es exactamente la topología de $Y$.
\end{enumerate}
\end{metodo}

\begin{ejemplo}[Ejemplo: problema IV.4]
$\tau_1=\{X,\vac,\s{a},\s{c,d},\s{a,c,d},\s{b,c,d,e}\}$ y $\tau_2=\{Y,\vac,\s{v},\s{w,z},\s{v,w,z},\s{w,x,y,z}\}$ tienen
abiertos de tamaños $0,1,2,3,4,5$ en ambos casos. Emparejando: $\s{a}\leftrightarrow\s{v}$ obliga $a\mapsto v$;
$\s{c,d}\leftrightarrow\s{w,z}$ obliga $\s{c,d}\mapsto\s{w,z}$; lo que queda, $\s{b,e}\mapsto\s{x,y}$.
Con $\varphi=(a\mapsto v,b\mapsto x,c\mapsto w,d\mapsto z,e\mapsto y)$ la tabla de imágenes recorre $\tau_2$:
$X\homeo Y$.

$\tau_3$ tiene 7 abiertos y las otras 6: $Z\not\homeo X$ y $Z\not\homeo Y$. Y todo espacio es homeomorfo a sí mismo.
\end{ejemplo}
\begin{ejemplo}[Ejemplo: problema IV.3(b)]
$X$ tiene 9 abiertos (10 con la corrección) y $Y$ tiene 5: no pueden ser homeomorfos, a pesar de lo que pide probar
el enunciado. Cuando un enunciado pide demostrar algo falso, la respuesta correcta es el contraargumento, bien
justificado.
\end{ejemplo}
\begin{ejemplo}[Ejemplos en $\R$]
$(-1,1)\homeo\R$ mediante $x\mapsto\dfrac{x}{1-|x|}$ (inversa $y\mapsto\dfrac{y}{1+|y|}$, ambas continuas). Dos
intervalos abiertos cualesquiera son homeomorfos (por una función afín). En cambio $[0,1]\not\homeo(0,1)$: la
compacidad es un invariante que se estudia más adelante en el curso.
\end{ejemplo}

\subsection{Clasificación de los espacios de tres puntos}
Las 29 topologías sobre $\s{a,b,c}$ se agrupan en \textbf{9 clases de homeomorfismo}\fuente{12,13}. Un representante de
cada una, con el número de topologías de la clase:
\begin{center}\small\renewcommand{\arraystretch}{1.2}
\begin{tabular}{@{}llc@{}}\toprule
representante & descripción & cuántas\\ \midrule
$\{\vac,X\}$ & indiscreta & 1\\
$\{\vac,\s{a},X\}$ & un punto abierto & 3\\
$\{\vac,\s{a,b},X\}$ & un par abierto & 3\\
$\{\vac,\s{a},\s{a,b},X\}$ & cadena & 6\\
$\{\vac,\s{a},\s{b,c},X\}$ & partición & 3\\
$\{\vac,\s{a},\s{b},\s{a,b},X\}$ & dos puntos abiertos y su unión & 3\\
$\{\vac,\s{a},\s{a,b},\s{a,c},X\}$ & un punto abierto y dos pares & 3\\
$\{\vac,\s{a},\s{b},\s{a,b},\s{a,c},X\}$ & dos puntos, dos pares & 6\\
$\Pt(X)$ & discreta & 1\\ \bottomrule
\end{tabular}\end{center}
$1+3+3+6+3+3+3+6+1=29$. Dos topologías de la misma fila se obtienen una de otra renombrando los puntos.
